% latexmk -pdf -pvc blatt10.tex
\documentclass{hott-übung}

\begin{document}

\setcounter{blattnummer}{10}
\setcounter{aufgabennummer}{0}

\blatt

Für Abbildungen \(f : B \to A\) und \(g : C \to A\) definieren wir den \emph{kanonischen Pullback}
\[B \underset{A}{\times} C \colonequiv (x : B) \times (y : C) \times (f(x) = g(y)).\]
Gegeben ein kommutatives Quadrat der Form
\begin{equation}
  \label{square}
  \begin{tikzcd}
  	D & B \\
  	C & A
  	\arrow["h", from=1-1, to=1-2]
  	\arrow["i"', from=1-1, to=2-1]
  	\arrow["S"{description}, draw=none, from=1-1, to=2-2]
  	\arrow["f", from=1-2, to=2-2]
  	\arrow["g"', from=2-1, to=2-2]
  \end{tikzcd}
\end{equation}
wobei \(S : f\circ h \sim g\circ i\), so erhalten wir eine kanonische Abbildung
\[\mathrm{gap} : D \longrightarrow B \underset{A}{\times} C, \, d \mapsto (h(d),i(d),S(d)).\]
Ein kommutatives Quadrat ist \emph{kartesisch}, falls diese Abbildung eine Äquivalenz ist.

\aufgabe{}
\begin{enumerate}
  \item
  Zeige, dass das folgende Quadrat (mit geeigneten Abbildungen) kartesisch ist.
  \[\begin{tikzcd}
    {(x \underset{A}{=} y)} & \eins \\
    \eins & A
    \arrow[from=1-1,to=1-2]
    \arrow[from=1-1,to=2-1]
    \arrow[from=1-2,to=2-2]
    \arrow[from=2-1,to=2-2]
  \end{tikzcd}
  \]
  \item
  Zeige: Eine Sequenz punktierter Abbildungen
  \[\begin{tikzcd}
    {\cdots} & {A_n} & {A_{n+1}} & {\cdots}
    \arrow[from=1-1,to=1-2,"f_{n-1}"]
    \arrow[from=1-2,to=1-3,"f_{n}"]
    \arrow[from=1-3,to=1-4,"f_{n+1}"]
  \end{tikzcd}\]
  ist genau dann eine Fasersequenz, wenn für alle \(n : \Z\), das folgende Quadrat kartesisch ist.
  \[\begin{tikzcd}
    {\fib_{f_n}(\ast)} & A_n \\
    \eins & A_{n+1}
    \arrow[from=1-1,to=1-2,"\pi_1"]
    \arrow[from=1-1,to=2-1]
    \arrow[from=1-2,to=2-2,"f"]
    \arrow[from=2-1,to=2-2]
  \end{tikzcd}\]
\end{enumerate}

\aufgabe{}
Zeige, dass das Quadrat (\ref{square}) genau dann kartesisch ist, wenn die Abbildung
\[\prod_{x : C} \fib_i(x) \to \fib_f(g(x)), \, (d,p) \mapsto (h(d),S(d)\kon\ap_g(p))\]
eine faserweise Äquivalenz ist.

\newpage
\aufgabe{Pullback-Pasting}
  Betrachte das kommutative Diagram
  \[\begin{tikzcd}
  	F & D & B \\
  	E & C & A
  	\arrow["k", from=1-1, to=1-2]
  	\arrow["l"', from=1-1, to=2-1]
  	\arrow["T"{description}, draw=none, from=1-1, to=2-2]
  	\arrow["h", from=1-2, to=1-3]
  	\arrow["i", from=1-2, to=2-2]
  	\arrow["S"{description}, draw=none, from=1-2, to=2-3]
  	\arrow["f", from=1-3, to=2-3]
  	\arrow["j"', from=2-1, to=2-2]
  	\arrow["g"', from=2-2, to=2-3]
  \end{tikzcd}\]
  mit Homotopien \(S : f\circ h \sim g\circ i\) und \(T : i\circ k \sim j \circ l\).
  Dann ist
  \[S_k \kon g(T) : f \circ h \circ k \sim g \circ j \circ l\]
  eine Homotopie, welche die Kommutativität des äußeren Rechtecks bezeugt.
  Unter der Voraussetzung, dass das rechte Quadrat kartesisch ist, zeige: das linke Quadrat ist genau dann kartesisch, wenn das äußere Rechteck kartesisch ist.

\aufgabe{}
  Sei \(f : B \to_\ast A\) eine punktierte Abbildung, sodass mit \(F \colonequiv \fib_f(\ast)\), die Sequenz \(F \xrightarrow{\pi_1} B \xrightarrow{f} A\) eine Fasersequenz wird.
  Verwende nun Pullback-Pasting, um eine Fasersequenz
  \[\begin{tikzcd}
    \Omega(B) & \Omega(A) & F & B & A
    \arrow[from=1-1,to=1-2]
    \arrow[from=1-2,to=1-3]
    \arrow[from=1-3,to=1-4]
    \arrow[from=1-4,to=1-5]
  \end{tikzcd}\]
  zu konstruieren.

% \aufgabe{}
% Zeige, dass die folgenden Quadrate mit geeigneten Abbildungen kartesisch sind.
% \begin{multicols}{2}
% \begin{enumerate}
%   \item
%   \begin{tikzcd}
%     {x = y} & \eins \\
%     \eins & A
%     \arrow[from=1-1,to=1-2]
%     \arrow[from=1-1,to=2-1]
%     \arrow[from=1-2,to=2-2]
%     \arrow[from=2-1,to=2-2]
%   \end{tikzcd}
%   für \(x,y : A\).
%   \item
%   \begin{tikzcd}
%     \fib_f(x) & B \\
%     \eins & A
%     \arrow[from=1-1,to=1-2]
%     \arrow[from=1-1,to=2-1]
%     \arrow[from=1-2,to=2-2,"f"]
%     \arrow[from=2-1,to=2-2]
%   \end{tikzcd}
%   für \(x : A\).
% \end{enumerate}
% \end{multicols}

\end{document}
